% theme: applications_of_speed
\MajorQuestionLesson{速さの利用問題}
\SetRowSpace{6mm}

\TypeQuestion{連立方程式をつくり、次の問いに答えなさい。}{deai}%大問の問題文はこの程度でOKです。

%各小問だけで問題を完結させてください。あと設定が似通いすぎているので「池の周りを回る」「運動場を一周する」など、ストーリー仕立てにしてください。
\QQ{1周1800mのコースを2人が同じ地点から同時に出発する。反対方向に進むと6分後に初めて出会い、同じ方向に進むと速い人が30分後に遅い人に初めて追いつく。}{}
\QQ{1周1680mの池の周りを2人が同じ地点から同時に出発する。反対方向に進むと6分後に初めて出会い、同じ方向に進むと速い人が21分後に遅い人に初めて追いつく。}{}
\QQ{1周1320mの公園の周回路を2人が同じ地点から同時に出発する。反対方向に進むと6分後に初めて出会い、同じ方向に進むと速い人が22分後に遅い人に初めて追いつく。}{}
\QQ{1周1440mのランニングコースを2人が同じ地点から同時に出発する。反対方向に進むと8分後に初めて出会い、同じ方向に進むと速い人が24分後に遅い人に初めて追いつく。}{}

\TypeQuestion{連立方程式をつくり、次の問いに答えなさい。}{deai}%大問の問題文はこの程度でOKです。

%ここも設定が似通いすぎている。数学的に意味のない設定でもいいので、状況がある程度リアルになるように、バリエーションをつけて。
\QQ{一定の速さで走る列車が、電柱の前を通過するのに12秒かかり、長さ180mのホームを通過するのに30秒かかった。この列車の速さと長さを求めなさい。}{}
\QQ{一定の速さで走る列車が、信号機の前を通過するのに8秒かかり、長さ240mの橋を渡り始めてから渡り終わるまでに20秒かかった。この列車の速さと長さを求めなさい。}{}
\QQ{一定の速さで走る列車が、長さ200mのトンネルに入り始めてから完全に出るまでに25秒、長さ500mのトンネルでは40秒かかった。この列車の速さと長さを求めなさい。}{}
\QQ{一定の速さで走る列車が、長さ90mのホームを通過するのに15秒、長さ210mの橋を渡り始めてから渡り終わるまでに25秒かかった。この列車の速さと長さを求めなさい。}{}
